
\chapter{Discussion}
\epigraph{``Another quote.''}{\vspace{10pt}Another Author}
\label{chapter:discussion}

\newpage

\section{CNN model and the regression problem}

\subsection{CNN model to identify cube regions and orientations}

To explore the task of locating the cube inside the measurement box, we followed the problem-complexity tiers described by \citeauthor{vapnikstatistical} in \cite{vapnikstatistical}. The model was first evaluated as a classifier to identify the cube's region and orientation. For the location-classification task across the acoustic sessions, the CNN achieved 98.59\% accuracy on the test set. This result indicates that the audio signals, once transformed into spectrograms, contain sufficient information for spatial localization and that the network can extract the features needed to identify the cube's region within the measurement box. The confusion matrix in Figure \ref{fig:cnn_confusion_matrix} shows a clear diagonal pattern, indicating strong classification performance. Twenty of the twenty-six classes achieved 100\% recall. Most errors occurred between neighboring classes. Class 1 had the most misclassifications, with 13 samples classified as class 0, resulting in 75\% accuracy for class 1. This confusion may be explained by the small distance between the points corresponding to these two sessions, as shown in Figure~\ref{fig:session_points}: S1 corresponds to class 0, and S2 corresponds to class 1.

In the orientation-classification task, the network was required to assign the cube's orientation to one of nine classes. The CNN achieved 29.06\% accuracy on the test set. Unlike the location-classification results, the confusion matrix in Figure \ref{fig:cnn_confusion_matrix_orientation} does not show a clear diagonal pattern. Instead, it reveals three clusters of misclassifications near the diagonal. In the first cluster, samples from classes 0, 1, and 2 are often classified as class 1. In the middle cluster, samples from classes 2, 3, 4, and 5 are often classified as class 4. In the rightmost cluster, samples from classes 5, 6, 7, and 8 are often classified as class 6. Overall, predictions are concentrated in classes 1 (10 degrees), 4 (40 degrees), and 6 (60 degrees). This pattern suggests that the network tends to favor a few broad orientation categories spaced approximately 20 to 30 degrees apart rather than reliably distinguishing all nine classes, which are spaced at 10-degree increments.

To further investigate whether the cube's proximity to the microphones and the orientation of its front face relative to the sides of the measurement box affected classification performance, the data were divided into three regions. The regional accuracy results in Table \ref{tab:regional_accuracy} show no improvement when the analysis is restricted to these regions. The lower accuracy may be explained by the reduced number of training samples in each region. Although accuracy remained low in this regional classification task, the confusion matrices in Figure \ref{fig:conf_matrix_3regions} show a more pronounced concentration of samples along the diagonal.

The cross-correlation plots in Figure~\ref{fig:cross_correlation_results} show that as the Euclidean distance between two locations increases, the normalized cross-correlation decreases, indicating that the acoustic signals become less similarity as the cube moves farther from the reference point. By contrast, changes in orientation do not produce a consistent decrease in cross-correlation; in some cases, it increases, indicating that the signals at certain orientations remain similar to the reference signal. This similarity may partly explain the CNN's poor performance on the orientation-classification task.

The two tasks show a pronounced difference in performance: spatial-location classification achieved over 98\% accuracy, whereas orientation classification achieved only 29.06\%. This difference may arise because rotating the cube changes the orientation of its surfaces without substantially changing their positions relative to the microphones and the wall in the box. Consequently, orientation may be encoded in subtle inter-channel phase shifts or in the attenuation of certain echoes, rather than in pronounced Time-Difference-Of-Arrival (TDOA) cues. By contrast, translating the cube changes the source–microphone geometry and may produce more distinct TDOA cues and changes in the reverberation pattern. The CNN may therefore have difficulty isolating orientation-dependent features in the spectrograms, which could explain why its predictions are concentrated in a few orientation classes.

\subsection{CNN model to estimate the cube location}

For regression, the CNN was used to estimate the cube's X, Y, and Z coordinates. Figures \ref{fig:five_fold_X_histo}, \ref{fig:five_fold_Y_histo}, and \ref{fig:five_fold_Z_histo} show the prediction-error histograms for the three coordinates across the five folds. The histograms indicate that the model generally localizes the cube with reasonable accuracy: most errors fall within approximately $\pm 2$ cm, a relatively narrow range compared with the dimensions of the measurement box and the cube. Most errors are concentrated around zero, suggesting that the network's predictions are close to the true positions. However, some larger errors occur, possibly because of challenging acoustic conditions, such as recording noise. The X and Y histograms show similar results. For Z, the distributions for folds 1 and 4 are broader, with more errors exceeding $\pm 2$ cm, suggesting that the network has greater difficulty estimating this coordinate. Overall, the histograms qualitatively demonstrate the CNN's strong spatial-localization capability, consistent with the high classification accuracy reported in the previous section. The error distributions are approximately bell-shaped and centered around zero, indicating good estimation performance across all three dimensions.

In this work, the MAE is the primary figure of merit used to evaluate a regression model's estimation performance \cite{friedmanrecomendation}. Table \ref{tab:mae_xyz_folds} summarizes the mean and standard deviation of the MAE for the X, Y, and Z coordinates across the five folds. The CNN achieves the lowest MAE for the Y coordinate, followed by the X coordinate. Although the X-coordinate MAE is slightly higher than the Y-coordinate MAE, the values are similar, at approximately 0.5 cm. For the Z coordinate, the mean MAE is nearly twice as high, at 0.9 cm, and the standard deviation is larger. These results are consistent with the histograms and further support the conclusion that the network estimates the cube's horizontal position more accurately than its vertical position.

Another metric used to evaluate regression performance is the Pearson correlation coefficient \cite{friedmanrecomendation}. Table \ref{tab:pearson_corr_xyz} reports this coefficient for all the folds and coordinates. The coefficients are consistently close to 1, indicating a strong linear relationship between the predicted and true coordinates. These results further support the network's ability to estimate the cube's position in the measurement box.

To analyze the error distribution across the specified ranges, Table \ref{tab:error_distribution_xyz} presents the percentage of samples whose estimation errors fall within each range. The results are consistent with the previous paragraph: Y is the most accurately estimated coordinate, with $92.90\%$ of samples falling within $\pm 1\text{ cm}$ and $98.32\%$ within $\pm 2\text{ cm}$, averaged across the five folds. X ranks second, with $87.69\%$ of samples falling within $\pm 1\text{ cm}$ and $97.83\%$ within $\pm 2\text{ cm}$, comparable to Y. Z has the largest fine-scale errors and the greatest variation across folds: only $64.44\%$ of samples fall within $\pm 1\text{ cm}$. The proportion within $\pm 3\text{ cm}$ rises to $98.32\%$, meaning that $8.92\%$ of Z-coordinate errors exceed $\pm 2\text{ cm}$. For all three coordinates, fewer than $0.6\%$ of samples have errors exceeding $\pm 4\text{ cm}$; equivalently, the proportions with errors within $\pm 4\text{ cm}$ are $99.42\%$ for X, $99.69\%$ for Y, and $99.45\%$ for Z.

The error distributions across the three dimensions are illustrated in Figure \ref{fig:plot3d_errors}, which shows the X, Y, and Z errors throughout the measurement box. The visualization indicates that errors are distributed relatively uniformly across the measurement space. Most errors are close to zero, while larger errors occur in a few regions; however, no clear clustering pattern is apparent. Figure \ref{fig:plot2d_errors} shows the error for each coordinate as a function of its corresponding coordinate value. For X, the absolute estimation error remains low near the center and increases toward the edges, where most outliers occur. A similar pattern is observed for Y, although the errors are lower than those observed for X. This suggests that the network estimates the Y coordinate more accurately than the X coordinate near the center. This result is notable because the network receives binaural input for X estimation, which could allow it to use the physical separation between the two microphones to learn interaural time differences. Despite this, Y estimation performs better. One possible explanation is that the cube's position along the Y axis produces distinctive interference patterns at different distances from the speaker and microphones. These changes in distance relative to the microphones may affect both the echo patterns and their magnitudes, particularly when the cube is near the speaker or the microphones. Such effects may help explain why Y estimation errors are lower than X estimation errors. The errors for Z are higher than those for X and Y. The Z-coordinate estimates may be affected by the uneven distribution of values around zero: the dataset contains more positive than negative values. Consequently, the network is trained on more positive samples than negative ones. This imbalance may contribute to the larger estimation errors for negative Z values. In addition, the robotic arm's base is located in a side (the base of the measurement box) perpendicular to the Z axis. Because this area is one of the two main surfaces contributing to binaural reverberation, it may complicate estimation in that part of the measurement space.


\section{ESN model and the regression problem}

As with the CNN, the ESN was first evaluated as a classifier to assess its ability to extract from the acoustic data the features needed to locate the cube within the measurement box. The ESN achieved 98.42\% accuracy, just 0.17 percentage points below the CNN's 98.59\%. The confusion matrix in Figure \ref{fig:confusion_ESN} shows a well-defined diagonal, indicating close agreement between the true and predicted labels. The few misclassifications are between neighboring classes. Unlike the CNN, the ESN also classifies sessions 0 and 1 more accurately; these sessions correspond to locations at the physical edge of the measurement area.

As with the CNN, the error distributions for the ESN estimators are shown in Figures \ref{fig:ESN_five_fold_X_histo}, \ref{fig:ESN_five_fold_Y_histo}, and \ref{fig:ESN_five_fold_Z_histo} for the X, Y, and Z coordinates, respectively. As in the CNN results, the histograms have approximately bell-shaped distributions centered around zero. However, the ESN error distributions are wider, indicating greater variability than that observed with the CNN.

The MAE values reported for the ESN in Table \ref{tab:ESN_mae_xyz_folds} are higher overall than those reported for the CNN in Table \ref{tab:mae_xyz_folds}. However, the relative performance of the X and Y estimators is reversed: the ESN achieves a lower MAE for X than for Y, while Z has the highest MAE of the three coordinates. The ESN's sensitivity to hyperparameter selection may partly explain these differences. Because the hyperparameters were initially tuned for regression on X, they may not have been optimal for Y, despite both coordinates using the same ground-truth transformation. The Z-coordinate data also have a different spatial distribution, and their ground-truth values do not require the arcsine normalization described in Equation \ref{eq:normalization}. These differences may further affect how well the selected hyperparameters generalize across ESN for diferent coordinates.

The error percentages reported across different tolerance ranges in Table \ref{tab:ESN_error_distribution_xyz} are consistent with the wider ESN error distributions shown in the histograms. A tolerance of $\pm 2$ cm captures a smaller proportion of the estimates; a tolerance of $\pm 3$ cm is needed to capture approximately 97\% of them.

The Pearson correlation coefficients on table \ref{tab:esn_pearson_corr_xyz} also exceed 0.99 in most cases. The only exception is the Z estimator in fold 3, suggesting that the regression model may have difficulty capturing the relationship between the features and target values for some samples in this fold.

Figure \ref{fig:esn_plot3d_errors} shows the spatial distribution of the estimation errors across all three dimensions. The scatter plot indicates that the error magnitudes are distributed relatively uniformly throughout the measurement region. However, the coordinate-wise plots in Figure \ref{fig:esn_plot2d_errors} show that the largest errors occur in a specific interior region rather than near the limits of the measurement space. This behavior may be explained by the multi-head architecture used for coordinate estimation. Two regression heads estimate the cube's position in the negative and positive coordinate ranges, respectively, while a third head selects which estimate to use. This selection may be less reliable near the center, where the two ranges meet. A fourth head identifies outliers by comparing the estimates; however, if the discrepancy between them is below the specified threshold, an incorrect head selection may not be flagged. This could explain why the central region remains prone to errors. Training both regression heads on an overlapping region around the boundary could mitigate this issue: each head would then learn from samples on both sides of the boundary, reducing the error when the selector chooses the less appropriate head.


\section{Beamforming}

\subsection{CNN with Beamforming}


\subsection{ESN with Beamforming}

\section{Model Comparison}
















%This indicates that while the ESN can still provide reasonably accurate estimates, it is more prone to larger errors compared to the CNN. Consequently, for applications where minimizing the largest errors is critical, the CNN may be the preferred choice despite its higher computational cost.

%Despite the higher MAE values, the ESN offers a significant reduction in model complexity compared to the CNN. The ESN requires only about 3\% of the trainable parameters used by the ResNet18 CNN, making it a more efficient option for the regression task with acoustic features as input. This efficiency could be particularly advantageous in scenarios with limited computational resources or when real-time processing is required. Overall, while the CNN achieves lower absolute errors, the ESN provides a compelling trade-off between accuracy and computational efficiency.
